\documentclass[11.5pt,a4paper]{article}
\usepackage[utf8]{inputenc}
\usepackage[T1]{fontenc}
\usepackage[french]{babel}
\usepackage{geometry}
\geometry{papersize={148mm,210mm}, top=2cm, bottom=2cm, left=1.5cm, right=1.5cm}

\usepackage{helvet}
\renewcommand{\familydefault}{\sfdefault}
\usepackage{xcolor}
\usepackage{tcolorbox}
\tcbuselibrary{skins, breakable}

\definecolor{CyrulnikBlue}{HTML}{2E4053}
\definecolor{HighlightGold}{HTML}{D4AC0D}

\newtcolorbox{FocusHumain}{
  colback=gray!5,
  colframe=CyrulnikBlue,
  fonttitle=\bfseries,
  title=Le Récit de la Résilience,
  arc=0mm,
  breakable
}

\begin{document}

\chapter*{Chapitre 1 : Le Cerveau Primitif et le Rapt de la Parole}

Le silence n'est pas une absence de bruit ; c'est parfois une cage. Dans le secret de nos boîtes crâniennes, un drame antique se joue chaque fois que nous étouffons une parole nécessaire. Pour comprendre pourquoi "dire non" est si douloureux pour certains, il faut descendre dans les caves de notre cerveau, là où la biologie commande la psychologie.

\section{L'Amygdale : Le Guetteur qui s'emballe}
Imaginez un petit noyau en forme d'amande, blotti au cœur de votre système limbique. C'est l'amygdale. Son métier est noble : elle assure votre survie. Elle est aux aguets, prête à hurler au moindre danger. Mais pour l'enfant à qui l'on a appris que déplaire est un crime, l'amygdale ne fait plus la différence entre un fauve et une simple contradiction. Chaque velléité d'affirmation déclenche un signal de détresse absolue.

\begin{FocusHumain}
\textit{L'histoire de Léo, 8 ans.} \\
Léo regarde ses chaussures. Son instituteur gronde, le ton monte, pas contre lui, mais l'ambiance électrique suffit. Léo sent son cœur battre dans sa gorge. Il voudrait expliquer qu'il n'a pas compris l'exercice, mais les mots sont comme collés au fond de sa poitrine. Son amygdale a pris le pouvoir : elle a coupé le courant vers ses centres de langage pour tout envoyer dans ses muscles de fuite... mais il ne peut pas fuir. Léo est "congelé". C'est le début du rapt de la parole.
\end{FocusHumain}

\section{Le "Hijack" émotionnel : Quand la pensée capitule}
Ce que les neurosciences appellent le \textit{Hijack}, ou détournement émotionnel, est un court-circuitage biologique. L'information sensorielle arrive à l'amygdale quelques millisecondes avant d'atteindre le Cortex Préfrontal — cette zone noble, située juste derrière votre front, qui nous permet de raisonner, de choisir nos mots, de poser des limites avec élégance.

Chez la personne non-assertive, cette "voie courte" devient l'autoroute principale. La peur gagne la course contre la raison. Avant même que vous ayez pu formuler une phrase calme ("Je ne suis pas d'accord avec toi"), votre corps a déjà décidé que vous étiez en danger de mort sociale. La parole est alors prise en otage par un système archaïque qui ne connaît que trois réponses : l'attaque (agressivité), la fuite (passivité) ou la sidération.

\section{Le vmPFC : La zone de la valeur de soi}
Grâce à l'Imagerie par Résonance Magnétique (IRM), nous observons aujourd'hui une zone fascinante : le Cortex Préfrontal Ventromédian (vmPFC). C'est ici que nous évaluons notre propre valeur. C'est ici que réside la force de dire "Je suis là, et mes besoins comptent". 

Les scans montrent que chez les individus souffrant d'une inhibition chronique, cette zone s'éteint lors des interactions sociales stressantes. Comme s'ils n'avaient plus neurologiquement le droit d'exister. À force de se taire, le cerveau finit par "penser" que nous sommes insignifiants.

\end{document}
